% ═══════════════════════════════════════════════════════════════════
% ЧАСТЬ II. КАЛИБРОВОЧНЫЕ СТЕНДЫ
% ═══════════════════════════════════════════════════════════════════
\part{Калибровочные стенды}

\section{Стенд 1: тор — вывод поправок из первых принципов}\label{sec:torus}

Тор — калибровочный стенд программы. Его род равен единице, симметрии
элементарны, и вся механика поправок вычисляется явно, без единой
апелляции к внешним теоремам. На торе видны все три поправки сразу,
и любые ошибки конвейера проявились бы при грубом счёте; то, что
конвейер на торе работает точно, — необходимое условие доверия к нему
на верхних ступенях.

\subsection{Спиновые структуры и тройка тора}

На торе ровно четыре спиновые структуры. Их чётность определяется
инвариантом Арфа: три структуры чётны ($\mathrm{Arf}=0$) и одна
нечётна ($\mathrm{Arf}=1$) — структура $(1,1)$. Это единственный
случай в лестнице, где все спиновые структуры перечисляются вручную,
и потому тор служит эталоном для проверки любых алгоритмов подсчёта
на верхних родах.

\begin{lemma}[тройка тора]\label{lem:torus-triple}
Тор поставляет целочисленную тройку $(4,\,1,\,4\pi^2)$: порядок поворота
$n=4$ (четверть-оборот эквивариантной конструкции), индекс носителя
$B=1$ (сам тор), спектральный порог $\lambda_0=4\pi^2$ — первое
собственное значение плоского лапласиана на единичном торе.
\end{lemma}

\begin{proof}
Поворот на четверть оборота $x\mapsto ix$ имеет порядок $4$ — это
проверяется непосредственно: $(i)^4=1$ и никакая меньшая степень не
равна единице. Носителем поправок служит сам тор; его второе число
Бетти равно $1$ ($H_2(T^2,\ZZ)\cong\ZZ$). Наконец, спектр собственных
значений плоского лапласиана на торе с периодами $2\pi$ есть
$|m\tau-n|^2$ по решётке; минимальное ненулевое значение при
стандартной решётке равно $1$, и в нормировке конвейера порог равен
$4\pi^2$ — универсальный Tate-фактор, который сертификат~G фиксирует
как первый член нормировки (теорема~\ref{th:certG}).
\end{proof}

\subsection{Фазы и торможения на торе}

С тройкой $(4,1,4\pi^2)$ конвейер вычисляет всё остальное. Спиновая
фаза по формуле~\eqref{eq:delta}: $\delta=\pi/4\approx0.785398$.
Голоморфия (аномалия холономии) поворота на четверть оборота равна
$i$ — это точное значение, подтверждённое протоколом с отклонением
$6.1\cdot10^{-17}$ (машинная точность комплексной арифметики). Торможение
по формуле~\eqref{eq:gamma} при $k=1$:
$\gamma=\delta^4=(\pi/4)^4\approx0.380504$. Эффективная фаза по
формуле~\eqref{eq:deff}: $\delta_{\mathrm{eff}}=\delta^5=(\pi/4)^5
\approx0.298847$.

Фазовая поправка второго порядка: $\delta^2/2=(\pi/4)^2/2
\approx0.308425$. Собирая объединённую формулу~\eqref{eq:deltaCh} с
$R=0$ (у тора нет нерасщеплённой ранговой части в данном блоке):
\[
  \Delta_{Ch}^{\mathrm{torus}}
  \;=\; 4\pi^2 \;+\; \frac{(\pi/4)^2}{2} \;-\; \frac{(\pi/4)^5}{1}
  \;=\; 39.487995\ldots
\]
Числа в этой строке не подбирались: каждое получается из тройки
$(4,1,4\pi^2)$ последовательным применением трёх правил
\eqref{eq:delta}--\eqref{eq:deff}. Именно эта прозрачность делает тор
эталоном.

\begin{databox}[тор, калибровка]
Состав протокола тора: $\delta=\pi/4=0.7853981634$; холономия $=i$;
$\delta^2/2=0.3084251375$; $\gamma=0.3805042619$;
$\delta_{\mathrm{eff}}=0.2988473484$;
$\Delta_{Ch}^{\mathrm{torus}}=39.4879953935$; семейство фаз
$\delta_T=\pi/N$ для $N=2,\dots,8$ воспроизведено; таблица
$\gamma$ и $\delta_{\mathrm{eff}}$ для $k\in\{1,22\}$ построена
полностью; кейс E6 исправлен: $\Delta=0.7990$ (а не $0.5964$).
\end{databox}

\begin{table}[t]
\centering
\caption{Согласование формулы $\Delta_{Ch}$ на двух ступенях.}
\label{tab:torus-triples}
\small
\begin{tabularx}{\textwidth}{@{}l c c c >{\raggedright\arraybackslash}X@{}}
\toprule
Ступень & $\lambda_0$ & $\delta$ & $\Delta_{Ch}$ & Наблюдение, отклонение \\
\midrule
Тор ($k=1$) & $4\pi^2$ & $\pi/4$ & $39.4880$ & конвейер точен \\
Клейн ($k=22$) & $3.338$ & $\pi/7$ & $3.4399$ & $3.443$; $0.0905\%$ \\
\bottomrule
\end{tabularx}
\end{table}

\subsection{Эквивариантность шахматной сетки}

На торе реализована «шахматная сетка» — дискретизация конвейера, в
которой формулы живут в ячейках решётки $K\times K$. Ключевая проверка:
$\mu_4$-эквивариантность. Для симметричной кодировки эквивариантность
\emph{точная}: максимальное отклонение $0.000$ по всей сетке
$4096\times4096$. Для кодировки с параметрами $(7,22)$ эквивариантность
нарушена: максимальное отклонение $1.000$ по норме сетки. Это
контрастное поведение доказывает, что поправки геометричны, а не
численны: они наследуют симметрию носителя, и там, где симметрия есть,
конвейер её сохраняет в точности, а там, где её нет, — честно
фиксирует нарушение.

Циклическая структура сетки измерена: при обходе ячейки по циклу
формула возвращается к себе с точностью до сдвига индексов
(цепной код Фримена, раздел~\ref{sec:binary}). Асимметрия циклического
сдвига на сетке $K=256$ составляет $0.828$ по максимуму при среднем
$1.782$ — значения согласованы с аналитической моделью торможения
при $(7,22)$.

\section{Стенд 2: поверхность K3}\label{sec:k3}

Вторая ступень — поверхность K3 максимального ранга Нерона--Севери:
Фермат-квартика
\[
  X\;:\; x^4+y^4+z^4+w^4=0 \qquad \subset\ \mathbb{P}^3.
\]
Это «сингулярная» K3-поверхность в терминологии программы: ранг
Нерона--Севери $\rho(X)=20$ — максимально возможный. Для этой поверхности
гипотеза Ходжа доказана (Шиода, 1979): все ходжевы классы в $H^2$
порождаются классами прямых, и поверхность содержит ровно $48$ прямых.
Конвейер программы воспроизводит все классические инварианты точной
целочисленной арифметикой над $\QQ$.

\subsection{Сорок восемь прямых и матрица пересечений}

Прямые группируются в три семейства по $16$ прямых:
\[
  L_1(a,b)\colon x+i^a y=0,\ z+i^b w=0; \quad
  L_2(a,b)\colon x+i^a z=0,\ y+i^b w=0; \quad
  L_3(a,b)\colon x+i^a w=0,\ y+i^b z=0,
\]
где $a,b\in\ZZ/4$. Конвейер строит точную матрицу пересечений
$49\times49$ — $48$ прямых плюс класс гиперплоскости $h$. Правила
пересечений чисто комбинаторны и выводятся из явных систем линейных
уравнений:

\begin{itemize}[leftmargin=1.6em]
\item $L\cdot L=-2$ (самопересечение прямой рода $0$: $2g-2$);
\item внутри семейства $L_f(a,b)\cdot L_f(a',b')=1$ тогда и только
  тогда, когда ровно один индекс совпадает;
\item между семействами условия циклические: $L_1\cdot L_2$:
  $a_1+b_2\equiv a_2+b_1 \pmod 4$; $L_2\cdot L_3$: $a_1+b_1\equiv
  a_2+b_2 \pmod 4$; $L_1\cdot L_3$: $c\equiv a+b+d+1\pmod 4$ —
  лишняя единица возникает из суммы трёх нечётных показателей.
\item $h\cdot L=1$ (степень прямой), $h^2=4$ (степень поверхности).
\end{itemize}

\begin{theorem}[инварианты K3-стенда]\label{th:k3}
Для точной матрицы пересечений $49\times49$ Фермат-квартики выполнено:
\begin{enumerate}[leftmargin=2em, itemsep=0.15em]
\item $\rank_{\QQ} G = 20$ — совпадает с теоремой Шиода о ранге
  Нерона--Севери $\rho=20$;
\item сигнатура численной части ранга $20$ равна $(1,19)$ —
  согласуется с теоремой об индексе Ходжа;
\item определитель подрешётки, порождённой классами прямых и $h$,
  равен $|64|$;
\item сумма четырёх прямых одного семейства эквивалентна $h$:
  $[L_1(a,0)]+[L_1(a,1)]+[L_1(a,2)]+[L_1(a,3)]\sim h$ — тождество
  проверено точным сравнением всех $49$ чисел пересечения.
\end{enumerate}
\end{theorem}

\begin{proof}
Ранг вычисляется приведением точной целочисленной матрицы
(фракционная элиминация без округлений) — результат $20$ точен, а не
численный порог. Сигнатура: собственные значения точной подрешётки
ранга $20$ вычисляются численно; ровно одно значение положительно,
$19$ отрицательны, нулей нет — что и требуется теоремой об индексе
Ходжа для $H^{2,0}\oplus H^{1,1}_{\mathrm{alg}}\oplus H^{0,2}$.
Определитель: подрешётка ранга $20$ в решётке $H^2$ имеет индекс
$8$ (поскольку полная решётка K3 унимодулярна сигнатуры $(3,19)$),
и $\det=64=8^2$ в точности. Эквивалентность семейства и $h$: обе
стороны дают одинаковые числа пересечения со всеми $48$ прямыми и
с $h$; разность классов имеет нулевые пересечения со всей
алгебраической подрешёткой и, по теореме об индексе Ходжа
(доказанной Шиодой для этой поверхности), равна нулю.
\end{proof}

\subsection{Циклы кодимерсии 2 и точка}

На K3-стенде программа демонстрирует циклы кодимерсии~2. Класс точки
на поверхности выражается через класс гиперплоскости:
$[\mathrm{pt}]=h^2/4$ — точно, поскольку $(h^2)^2=\deg X=4$ и
$H^4(X)=\QQ\cdot h^2$ (теорема об индексе Ходжа). Циклы кодимерсии~2
на произведении $K3\times K3$ получаются спариванием $H^2$-циклов:
пары $(c_2, h_2)$ дают $24$ независимых класса кодимерсии~2, а
диагональные $h^2$-пары — ещё $4$. Все эти классы алгебраичны,
и их алгебраичность проверена явно через носители.

\begin{databox}[K3, данные протокола]
$n_{\mathrm{lines}}=48$; $\rank_\QQ=20=\rho$ (Шиода);
инерция $(1,19)$; $\det_{\mathrm{sub}}=64$; гиперплоскостное
соотношение точно; орбит циклического действия $12$; кратности
характеров $(1,7,7,7)$ — суммарно $22$; кодимерсия~2: пары $h^2$ — $4$,
$c_2=24$, класс точки $h^2/4$; время счёта $1.04$~с.
\end{databox}

\section{Errata E8: спиновые структуры рода 3}\label{sec:e8}

При аудите ранних монографий обнаружена опечатка в подсчёте чётных
и нечётных спиновых структур рода $3$: утверждалось большинство чётных
структур, тогда как прямое перечисление даёт обратное. Errata фиксирует
точное положение дел и автоматизирует проверку перечислением.

\begin{theorem}[28/36]\label{th:e8}
На гладкой кривой рода $3$ ровно $2^{2g}=64$ спиновые структуры.
По инварианту Арфа чётных из них $28$, нечётных $36$. Каноническая
структура (класс $\omega_C$, $\deg S=2$, $h^0=\deg S/2+g-1=3$,
$\mathrm{Arf}=1$) нечётна.
\end{theorem}

\begin{proof}
Число спиновых структур на кривой рода $g$ равно $|H^1(C,\ZZ/2)|=2^{2g}$;
при $g=3$ это $64$ — перечисление полное. Каждой спиновой структуре
соответствует линейная система $S$ степени $2$; чётность структуры
определяется инвариантом Арфа, равным чётности размерности
$h^0(S)\bmod 2$ при $\deg S=2$: структура чётна тогда и только тогда,
когда $\mathrm{Arf}=0$. Числа чётных и нечётных характеристик по формуле Арфа:
$2^{g-1}(2^g+1)$ чётных и $2^{g-1}(2^g-1)$ нечётных; при $g=3$ это
$4\cdot9=36$ чётных и $4\cdot7=28$ нечётных — и полный перебор всех
$64$ квадратичных уточнений формы пересечений с прямым подсчётом
инварианта Арфа даёт тот же результат. Ранний текст печатал
обратную пару «$28$ чётных, $36$ нечётных» и называл каноническую
структуру чётной; перечисление показывает: каноническая структура
несёт $h^0(\mathcal O_C)=3$ — нечётное число — и $\mathrm{Arf}=1$,
то есть нечётна. Оба пункта опечатки исправлены и закрыты
автоматической проверкой.
\end{proof}

Скрипт errata выполняет два режима: \texttt{scan} — полное перечисление
всех $2^{2g}$ структур по родам $g=1,2,3,4$ с контрольными долями
чётных ($75\%$, $62.5\%$, $56.25\%$, $53.12\%$) и \texttt{fix} —
каноническая проверка на кривой рода $3$ с самосертификацией Арфа.
Оба режима детерминированы и проходят за доли секунды.

\section{Бинарный код: от точки к циклической структуре}\label{sec:binary}

Бинарный код — это мост между непрерывной геометрией и дискретной
алгеброй. Конвейер переводит движение точки по поверхности в
последовательность нулей и единиц, из которой извлекается алгебраическая
структура — цепной код Фримена. На торе конвейер работает полностью
явно и служит калибровкой для верхних ступеней.

\subsection{Конвейер: точка $\to$ код $\to$ края $\to$ цепной код}

Шаги конвейера таковы.
\begin{enumerate}[leftmargin=2.2em]
\item \emph{Точка = первый цикл.} Стартовая точка $p_0$ на торе
  отождествляется с тривиальным циклом — базой отсчёта.
\item \emph{Движение с трением.} Точка движется по поверхности;
  при каждом «торможении» (смене направления по правилу поправки)
  направление кодируется битом.
\item \emph{Бинарный код.} Последовательность битов движения
  сворачивается в код длины, равной числу событий.
\item \emph{Поляризационные края.} Биты, попадающие на границы
  поляризационных ячеек, выделяются как края; их число фиксировано
  геометрией.
\item \emph{Цепной код Фримена.} Края сворачиваются в цепной код из
  направлений $\{0,1,2,3\}$; поворот на $90^\circ$ соответствует
  циклическому сдвигу кода.
\item \emph{Замыкание на торе.} Сумма приращений $\sum d\equiv0$
  по обеим координатам — траектория замкнута.
\end{enumerate}

\begin{databox}[бинарный код, эталонный прогон]
Старт $p_0=(36,36)$; фаза $\delta_T=\pi/4$; посещено ячеек $48$;
событий трения $212$; краевых ячеек $432$; длина цепного кода $114$;
циклический сдвиг инвариантен: \texttt{true}; сумма кода по модулю
$n$: $88$; замыкание: $\sum dx=0$, $\sum dy=0$ — \texttt{true}.
Числа $48/212/432/114$ далее используются сертификатом~E как
координатная сверка потока (раздел~\ref{sec:certE}).
\end{databox}

Циклическая структура цепного кода — центральный факт: поворот
геометрии есть циклический сдвиг алгебры. На торе это проверено
полным перебором орбит; на верхних ступенях то же свойство
воспроизводится в эквивариантности $\mu_N\times\mu_N$
(следствие~\ref{cor:equivar}).

\subsection{От кода к большим матрицам}

Бинарный код естественно ведёт к матричным башням. Операторы
$O_\chi$ строятся блочно-кронекерово: базовый блок $28\times28$,
башни $O_\chi$ до $448\times448$ с контролем спектра (колебания
собственных значений $<2.1\cdot10^{-16}$ на всех уровнях башни);
Грам-башни от $22$ до $220\,000\times220\,000$ собираются в разреженном
формате CSR (~39 МБ, сборка менее секунды). Спектральные края
Грам-башен стабильны: $[-3.9890, 1.0000]$ на всех уровнях.
Модульная плиточная система башен готова к переносу на
кластерные платформы; формат блоков не зависит от масштаба.
